% ============================================================================
% DRAFT v2 -- Section 6: Honest positioning
% Published values are as reported by the original authors (not re-run); the
% controlled comparison is our own (proxy_diag/phaseB_truth.json, 210 detectors).
% ============================================================================

\section{Discussion: Where This Pipeline Stands}
\label{sec:positioning}

This paper reports a diagnostic study, not a new state-of-the-art detector, and it
is therefore necessary to state plainly where the pipeline of
Section~\ref{sec:method} stands. Two comparisons are relevant. As \emph{reported by
the original authors}, image-level AUROC on MVTec-AD runs from $78.0\%$
(AutoEncoder) and $85.0\%$ (VAE)~\cite{bergmann2018improving} through $95.6\%$
(STPM)~\cite{wang2021stpm}, $96.0\%$ (SPADE)~\cite{cohen2020spade}, $97.9\%$
(PaDiM)~\cite{defard2021padim}, $98.5\%$
(reverse distillation)~\cite{deng2022reverse}, $99.1\%$
(PatchCore~\cite{roth2022patchcore}, EfficientAD~\cite{batzner2024efficientad}) and
$99.3$--$99.6\%$
(CFA~\cite{lee2022cfa}, FastFlow~\cite{yu2021fastflow},
SimpleNet~\cite{liu2023simplenet}, Dinomaly~\cite{guo2025dinomaly}). FastFlow
reports $99.4\%$. These numbers were not re-run here and their protocols are not
identical to ours, so they are for orientation only. Two of them, SPADE and
FastFlow, are preprints that were never published at a venue; they are included
because they are standard baselines in this literature, and their preprint status
is declared here rather than left implicit. The second comparison,
Table~\ref{tab:controlled}, is our own, in which all 14 backbones of the search
space were trained under one protocol on all 15 categories.

\begin{table}[t]
\centering
\caption{Controlled comparison: every backbone of the search space trained on
every MVTec-AD category under an identical protocol (500 epochs, batch size 64,
seed 42, ResNet-18 teacher, reverse distillation; 210 detectors). Deployed size is
the student encoder plus decoder. The manuscript's selected architecture is marked
with~$\dagger$. Mean AUROC and Rank are single-seed (seed 42) and should be read
with the resolution of Section~\ref{subsec:resolvability} in mind: \emph{every}
adjacent
rank differs by less than the $2.98$-point minimum detectable difference (the
largest adjacent gap is $0.75$ points), and only non-adjacent extremes --- for
instance rank $1$ against rank $13$, a gap of $3.23$ points --- exceed it. The
ordering is therefore resolved only at its extremes and unresolved across most of
its length. Six of the fourteen backbones were replicated at two further seeds; among
those six the ordering is \emph{identical} in all three seeds
(\texttt{resnet18} $>$ \texttt{efficientnet\_b0} $>$
\texttt{mobilenetv3\_small\_100} $>$ ViT-Tiny $>$ ConvNeXtV2-Tiny $>$
Swin-Tiny, maximum rank movement $0$), which is the evidence that the axis is
real but too small to be ranked by a single-seed comparison of the other eight.}
\label{tab:controlled}
\resizebox{\columnwidth}{!}{%
\begin{tabular}{lrrr}
\toprule
Backbone & Student (M) & Mean AUROC (\%) & Rank \\
\midrule
ResNet-18        & 13.04 & \textbf{83.34} & 1 \\
ResNet-34        & 23.15 & 82.99 & 2 \\
ResNet-50        & 25.77 & 82.75 & 3 \\
MobileOne-S0     &  6.27 & 82.58 & 4 \\
EfficientNet-B0$^\dagger$ & 5.41 & 82.47 & 5 \\
DeiT-Tiny        &  7.31 & 82.03 & 6 \\
MobileNetV3-Small & 2.81 & 81.85 & 7 \\
ViT-Tiny         &  7.31 & 81.51 & 8 \\
ConvNeXt-Small   & 51.39 & 81.43 & 9 \\
ConvNeXtV2-Tiny  & 29.80 & 81.35 & 10 \\
ConvNeXt-Tiny    & 29.75 & 80.91 & 11 \\
EfficientNet-B1  &  7.92 & 80.43 & 12 \\
Swin-Tiny        & 29.26 & 80.11 & 13 \\
ConvNeXtV2-Base  & 89.69 & 79.36 & 14 \\
\bottomrule
\end{tabular}}

\end{table}

The pipeline described in this paper reaches $82.47\%$ mean AUROC on MVTec-AD with
the architecture the search selected, \texttt{efficientnet\_b0} (5.41M parameters),
against reported values of $95.6\%$--$99.6\%$ for the established detectors. We
therefore do not claim a competitive operating point: the gap to the strongest
published single-class detectors is $17.1$ points, and we report it
rather than obscure it. We can also name at least one candidate cause. The
deployed detector uses an ImageNet-1k ResNet-18 as its frozen teacher
(Section~\ref{subsec:teacher}), whereas reverse distillation as published
pre-trains the teacher on the normal images of the target category; the same code
path also fixes the student bottleneck at a width the search cannot change. We do
not claim that these account for the whole gap, and we have not run the
teacher-adapted variant, so the attribution is a hypothesis rather than a result.
A second caveat is that the harness itself was not validated against a published
reverse-distillation configuration: we did not reproduce a reference RD number
inside our code path, so the $17.1$-point gap could in principle contain a
harness-level difference (for instance the batch-norm behaviour of the largest
backbones, which are trained with gradient accumulation and a smaller micro-batch;
Section~\ref{subsec:implementation}) on top of the teacher choice. We report this
rather than present the comparison as exact, and we note that it does not affect
the paper's diagnostic claims, which are all within-harness. Nor is the gap
explained by the parameter budget:
Section~\ref{subsec:capacity_null} varies the deployed size over a $9.2\times$
range across five capacity bands and finds no statistically resolvable effect, and
across the full space of Table~\ref{tab:controlled} the deployed size spans
$32\times$ while the backbone accounts for $0.54\%$ of the variance. One caveat
applies to the reported values: they come from implementations, augmentation
policies and training schedules that differ from ours, so the comparison is
indicative rather than exact. Table~\ref{tab:controlled} is the within-study
comparison on which the analysis in Section~\ref{sec:diagnosis} rests; it was
produced by the same code path, on the same data split, with one variable, the
backbone, changed at a time.

Two things are worth taking from Table~\ref{tab:controlled} in their own right.
It is, to our knowledge, the first comparison of a full anomaly-detection search
space under a single protocol, and as such it is a reusable resource: the ordering
it establishes is what makes the negative result in Section~\ref{sec:diagnosis}
measurable rather than anecdotal. And it shows that the search's output, at rank 5
of 14, is a reasonable but unremarkable choice from a space whose top eleven
candidates span $2.43$ points. The canonical headline number of this paper is the
$82.47\%$ mean AUROC of \texttt{efficientnet\_b0} on the fifteen categories, read
from the seed-42 truth matrix
(\texttt{proxy\_diag/phaseB\_truth.json}). An earlier, independently trained run of
the same configuration obtained $82.40\%$ on the same fifteen categories, a
difference of $0.07$ points. We report that agreement as a stability check and not
as confirmation: $0.07$ points is far below this paper's own $2.98$-point minimum
detectable difference, so the two numbers are simply \emph{consistent with} a
stable protocol, and the sentence cannot be made stronger without contradicting
Section~\ref{subsec:resolvability}. The search itself does \emph{not} reproduce as
closely;
Section~\ref{subsec:proxy_stability} documents why.
